% GENERATED FILE. Edit canonical lessons, then run npm run generate:books.
% canonical-source-hash: fnv1a64:6a7986b3025ebe57
% canonical-lessons: MR-C210-dantavaidya, MR-C210-khatik, MR-C210-karagir, MR-C210-kamgar, MR-C210-karmacari, MR-R210-first-pass-words-from-chapters-202-205, MR-R210-second-pass-words-from-chapters-206-210

\chapter{Dentist, Butcher, Craftsman, Worker, Employee}
\label{ch:mr-dentist-butcher-craftsman-worker-employee}
\hlchaptermodality{210}

\begin{chapteropening}
\textbf{By the end of this chapter:} \emph{I can say a dentist, a butcher, a craftsman, a worker and an employee.}

Complete the last lesson of chapter 210: I can say a dentist, a butcher, a craftsman, a worker and an employee.
\end{chapteropening}

\section[\texorpdfstring{\mr{दंतवैद्य} (dantavaidya)}{दंतवैद्य (dantavaidya)}]{\mr{दंतवैद्य} (dantavaidya) — a dentist}

\label{lesson:MR-C210-dantavaidya}



\begin{quote}
Before the new one: say the Marathi for an army, then the Marathi for a soldier.
\end{quote}

\subsection*{You'll want to know: \mr{दंतवैद्य}}
\textbf{\mr{दंतवैद्य}} — \emph{dantavaidya} — \textquotedblleft{}a dentist\textquotedblright{}.

\begin{grammarlens}[title={What you've built}]
One more word for this chapter.
\end{grammarlens}

\subsection*{Your turn}
\begin{itemize}
\raggedright
  \item \emph{Say it:} \emph{dantavaidya}
  \item \emph{Say it:} \emph{dantavaidya}, once more
  \item \emph{Say it:} say \emph{dantavaidya}
  \item \emph{Recall:} say the Marathi for an army, then the Marathi for a soldier, then say \emph{dantavaidya} again
\end{itemize}

\begin{tcolorbox}[breakable,colback=teal!4,colframe=teal!35!black,title={Before you move on}]
What does \mr{दंतवैद्य} mean? (\textquotedblleft{}A dentist\textquotedblright{}.) Say it once more.
\end{tcolorbox}
\section[\texorpdfstring{\mr{खाटीक} (khāṭīk)}{खाटीक (khāṭīk)}]{\mr{खाटीक} (khāṭīk) — a butcher}

\label{lesson:MR-C210-khatik}



\begin{quote}
Before the new one: say the Marathi for a soldier, then the Marathi for a dentist.
\end{quote}

\subsection*{You'll want to know: \mr{खाटीक}}
\textbf{\mr{खाटीक}} — \emph{khāṭīk} — \textquotedblleft{}a butcher\textquotedblright{}.

\begin{grammarlens}[title={What you've built}]
One more word for this chapter.
\end{grammarlens}

\subsection*{Your turn}
\begin{itemize}
\raggedright
  \item \emph{Say it:} \emph{khāṭīk}
  \item \emph{Say it:} \emph{khāṭīk}, once more
  \item \emph{Say it:} say \emph{khāṭīk}
  \item \emph{Recall:} say the Marathi for a soldier, then the Marathi for a dentist, then say \emph{khāṭīk} again
\end{itemize}

\begin{tcolorbox}[breakable,colback=teal!4,colframe=teal!35!black,title={Before you move on}]
What does \mr{खाटीक} mean? (\textquotedblleft{}A butcher\textquotedblright{}.) Say it once more.
\end{tcolorbox}
\section[\texorpdfstring{\mr{कारागीर} (kārāgīr)}{कारागीर (kārāgīr)}]{\mr{कारागीर} (kārāgīr) — a craftsman}

\label{lesson:MR-C210-karagir}



\begin{quote}
Before the new one: say the Marathi for a dentist, then the Marathi for a butcher.
\end{quote}

\subsection*{You'll want to know: \mr{कारागीर}}
\textbf{\mr{कारागीर}} — \emph{kārāgīr} — \textquotedblleft{}a craftsman\textquotedblright{}.

\begin{grammarlens}[title={What you've built}]
One more word for this chapter.
\end{grammarlens}

\subsection*{Your turn}
\begin{itemize}
\raggedright
  \item \emph{Say it:} \emph{kārāgīr}
  \item \emph{Say it:} \emph{kārāgīr}, once more
  \item \emph{Say it:} say \emph{kārāgīr}
  \item \emph{Recall:} say the Marathi for a dentist, then the Marathi for a butcher, then say \emph{kārāgīr} again
\end{itemize}

\begin{tcolorbox}[breakable,colback=teal!4,colframe=teal!35!black,title={Before you move on}]
What does \mr{कारागीर} mean? (\textquotedblleft{}A craftsman\textquotedblright{}.) Say it once more.
\end{tcolorbox}
\section[\texorpdfstring{\mr{कामगार} (kāmgār)}{कामगार (kāmgār)}]{\mr{कामगार} (kāmgār) — a worker}

\label{lesson:MR-C210-kamgar}



\begin{quote}
Before the new one: say the Marathi for a butcher, then the Marathi for a craftsman.
\end{quote}

\subsection*{You'll want to know: \mr{कामगार}}
\textbf{\mr{कामगार}} — \emph{kāmgār} — \textquotedblleft{}a worker\textquotedblright{}.

\begin{grammarlens}[title={What you've built}]
One more word for this chapter.
\end{grammarlens}

\subsection*{Your turn}
\begin{itemize}
\raggedright
  \item \emph{Say it:} \emph{kāmgār}
  \item \emph{Say it:} \emph{kāmgār}, once more
  \item \emph{Say it:} say \emph{kāmgār}
  \item \emph{Recall:} say the Marathi for a butcher, then the Marathi for a craftsman, then say \emph{kāmgār} again
\end{itemize}

\begin{tcolorbox}[breakable,colback=teal!4,colframe=teal!35!black,title={Before you move on}]
What does \mr{कामगार} mean? (\textquotedblleft{}A worker\textquotedblright{}.) Say it once more.
\end{tcolorbox}
\section[\texorpdfstring{\mr{कर्मचारी} (karmacārī)}{कर्मचारी (karmacārī)}]{\mr{कर्मचारी} (karmacārī) — an employee}

\label{lesson:MR-C210-karmacari}



\begin{quote}
Before the new one: say the Marathi for a craftsman, then the Marathi for a worker.
\end{quote}

\subsection*{You'll want to know: \mr{कर्मचारी}}
\textbf{\mr{कर्मचारी}} — \emph{karmacārī} — \textquotedblleft{}an employee\textquotedblright{}.

\begin{grammarlens}[title={What you've built}]
One more word for this chapter.
\end{grammarlens}

\subsection*{Your turn}
\begin{itemize}
\raggedright
  \item \emph{Say it:} \emph{karmacārī}
  \item \emph{Say it:} \emph{karmacārī}, once more
  \item \emph{Say it:} say \emph{karmacārī}
  \item \emph{Recall:} say the Marathi for a craftsman, then the Marathi for a worker, then say \emph{karmacārī} again
\end{itemize}

\begin{tcolorbox}[breakable,colback=teal!4,colframe=teal!35!black,title={Before you move on}]
What does \mr{कर्मचारी} mean? (\textquotedblleft{}An employee\textquotedblright{}.) Say it once more.
\end{tcolorbox}
\section[(dialogue)]{Words from chapters 202-205 — a first pass}

\label{lesson:MR-R210-first-pass-words-from-chapters-202-205}



\begin{quote}
Say the words for a worker and an employee.
\end{quote}

\subsection*{Your turn}
Say these aloud:

\begin{itemize}
\raggedright
  \item a border, a region, a capital city, a suburb, up
  \item certainly, really, almost, about, only
  \item serious, shy, valiant, fearful, curious
  \item free of charge, modern, ancient, cooked, ripe
\end{itemize}

\begin{tcolorbox}[breakable,colback=teal!4,colframe=teal!35!black,title={Before you move on}]
A border? (\textbf{\mr{सीमा}}, \emph{sīmā}.) A region? (\textbf{\mr{प्रदेश}}, \emph{pradeś}.) A capital city? (\textbf{\mr{राजधानी}}, \emph{rājdhānī}.) A suburb? (\textbf{\mr{उपनगर}}, \emph{upnagar}.) Up? (\textbf{\mr{वरती}}, \emph{varatī}.) Certainly? (\textbf{\mr{नक्की}}, \emph{nakkī}.) Really? (\textbf{\mr{खरोखर}}, \emph{kharokhar}.) Almost? (\textbf{\mr{जवळजवळ}}, \emph{javaḷjavaḷ}.) About? (\textbf{\mr{सुमारे}}, \emph{sumāre}.) Only? (\textbf{\mr{फक्त}}, \emph{phakta}.) Serious? (\textbf{\mr{गंभीर}}, \emph{gambhīr}.) Shy? (\textbf{\mr{लाजाळू}}, \emph{lājāḷū}.) Valiant? (\textbf{\mr{शूर}}, \emph{śūr}.) Fearful? (\textbf{\mr{भित्रा} / \mr{भित्री}}, \emph{bhitrā / bhitrī}.) Curious? (\textbf{\mr{उत्सुक}}, \emph{utsuk}.) Free of charge? (\textbf{\mr{मोफत}}, \emph{mophat}.) Modern? (\textbf{\mr{आधुनिक}}, \emph{ādhunik}.) Ancient? (\textbf{\mr{प्राचीन}}, \emph{prācīn}.) Cooked? (\textbf{\mr{शिजलेला}}, \emph{śijlelā}.) Ripe? (\textbf{\mr{पिकलेला}}, \emph{piklelā}.)
\end{tcolorbox}
\section[(dialogue)]{Words from chapters 206-210 — a second pass}

\label{lesson:MR-R210-second-pass-words-from-chapters-206-210}



\begin{quote}
Say the words for public and private.
\end{quote}

\subsection*{Your turn}
Say these aloud:

\begin{itemize}
\raggedright
  \item public, private, personal, ordinary, special
  \item bold, careful, careless, busy, free
  \item an armchair, a button, a snack, a tip; a prize, a stop
  \item the economy, war, peace, an army, a soldier
  \item a dentist, a butcher, a craftsman, a worker, an employee
\end{itemize}

\begin{tcolorbox}[breakable,colback=teal!4,colframe=teal!35!black,title={Before you move on}]
Public? (\textbf{\mr{सार्वजनिक}}, \emph{sārvajanik}.) Private? (\textbf{\mr{खाजगी}}, \emph{khājgī}.) Personal? (\textbf{\mr{वैयक्तिक}}, \emph{vaiyaktik}.) Ordinary? (\textbf{\mr{सामान्य}}, \emph{sāmānya}.) Special? (\textbf{\mr{खास}}, \emph{khās}.) Bold? (\textbf{\mr{धीट}}, \emph{dhīṭ}.) Careful? (\textbf{\mr{सावध}}, \emph{sāvadh}.) Careless? (\textbf{\mr{निष्काळजी}}, \emph{niṣkāḷjī}.) Busy? (\textbf{\mr{व्यस्त}}, \emph{vyasta}.) Free? (\textbf{\mr{मोकळा}}, \emph{mokḷā}.) An armchair? (\textbf{\mr{आरामखुर्ची}}, \emph{ārāmkhurcī}.) A button? (\textbf{\mr{बटण}}, \emph{baṭaṇ}.) A snack? (\textbf{\mr{खाऊ}}, \emph{khāū}.) A tip; a prize? (\textbf{\mr{बक्षीस}}, \emph{bakṣīs}.) A stop? (\textbf{\mr{थांबा}}, \emph{thāmbā}.) The economy? (\textbf{\mr{अर्थव्यवस्था}}, \emph{arthavyavasthā}.) War? (\textbf{\mr{युद्ध}}, \emph{yuddha}.) Peace? (\textbf{\mr{शांतता}}, \emph{śāntatā}.) An army? (\textbf{\mr{सैन्य}}, \emph{sainya}.) A soldier? (\textbf{\mr{सैनिक}}, \emph{sainik}.) A dentist? (\textbf{\mr{दंतवैद्य}}, \emph{dantavaidya}.) A butcher? (\textbf{\mr{खाटीक}}, \emph{khāṭīk}.) A craftsman? (\textbf{\mr{कारागीर}}, \emph{kārāgīr}.) A worker? (\textbf{\mr{कामगार}}, \emph{kāmgār}.) An employee? (\textbf{\mr{कर्मचारी}}, \emph{karmacārī}.)
\end{tcolorbox}
